% Part: first-order-logic
% Chapter: model-theory
% Section: substructures

\documentclass[../../../include/open-logic-section]{subfiles}

\begin{document}

\olfileid{mod}{bas}{sub}
\olsection{子\printtoken{p}{structure}}

一个结构~$\Struct{M}$ 的论域可能是另一个结构~$\Struct{M'}$ 的子集。
但显然，只有当 $\Domain{M} \subseteq \Domain{M'}$，
且 $\Struct{M}$ 与 $\Struct{M'}$ 对语言符号的解释至少在共享部分~$\Domain{M}$ 上“一致”时，
我们才应将 $\Struct{M}$ 视作 $\Struct{M'}$ 的“部分”。

\begin{defn}
\ollabel{defn:substructure}
给定同一语言~$\Lang L$ 的两个结构 $\Struct M$ 和 $\Struct M'$，
我们称 $\Struct M$ 是 $\Struct M'$ 的\emph{子结构}，
而称 $\Struct M'$ 是 $\Struct M$ 的\emph{扩张}，
记作 $\Struct M \substruct \Struct M'$，当且仅当
\begin{enumerate}
\item $\Domain{M} \subseteq \Domain{M'}$，
\item 对每个常项符号 $c \in \Lang L$，$\Assign{c}{M} =
    \Assign{c}{M'}$；
\item 对每个 $n$ 元函数符号 $f \in \Lang L$
  $\Assign{f}{M}(a_1, \dots, a_n) = \Assign{f}{M'}(a_1, \dots, a_n)$
  对所有 $a_1$、\ldots、$a_n \in \Domain{M}$ 成立。
\item 对每个 $n$ 元谓词符号 $R \in \Lang L$，$\langle
  a_1, \dots, a_n\rangle \in \Assign{R}{M}$ 当且仅当 $\langle a_1, \dots,
  a_n\rangle \in \Assign{R}{M'}$ 对所有 $a_1$、\ldots、$a_n \in
  \Domain{M}$ 成立。
\end{enumerate}
\end{defn}

\begin{rem}
\ollabel{rem:substructure}
如果语言中没有常项符号或函数符号，那么对任意 $N
\subseteq \Domain{M}$，令 $\Assign{R}{N} =
\Assign{R}{M} \cap N^n$，就得到一个以 $N$ 为论域（$\Domain{N} = N$）的 $\Struct M$ 的子结构~$\Struct{N}$。
\end{rem}

% prove something about this? Examples?

\end{document}